\section{Analysis}
\label{sec:analysis}

In this section we study the formal properties of \CBSOMminus\ under
cognitive bias functions from the Receptive-Resistant region $R$.
We first establish the structural invariants that persist from \SOMminus\
to \CBSOMminus\ (Section~\ref{sec:analysis-structure}), then introduce
the \emph{scaled linear bias family} $\mathcal{F}_\alpha$
(Section~\ref{sec:scaled-linear}), and finally prove consensus
guarantees for clique graphs with $n\ge 2$
(Section~\ref{sec:consensus-cliques}).

% ===================================================================
\subsection{Structural Properties under \texorpdfstring{$R$}{R}-Biases}
\label{sec:analysis-structure}

A central motivation for restricting to the Receptive-Resistant region
$R$ is that continuous $R$-biases guarantee consensus in the synchronous
DeGroot setting on strongly connected graphs~\cite{Alvim2024}. We ask
whether the silence mechanism of \CBSOMminus\ preserves the structural
invariants that underpin the \SOMminus\ consensus proofs~\cite{aranda2024sound}. We show that the answer is affirmative: the silence criterion in
Eq.~\eqref{eq:cbsom-minus-s} depends solely on raw opinion proximity
$|B_i^t-B_j^t|\le\tau_i$---not on the bias-filtered update
value---so the cognitive bias layer \emph{decouples} from the silence
decision.

\begin{lemma}[All-silent state lasts at most one round]
\label{lem:no-perpetual-silence}
  Let $(G,\mathbf{B}^0,\mathbf{S}^0,\mu_G)$ be a \CBSOMminus\ model.
  For any $t\in\mathbb{N}$, if $S_i^t=0$ for all $i\in A$, then
  $S_i^{t+1}=1$ for all $i\in A$.
\end{lemma}

\begin{proof}
  If $S_i^t=0$ for all $i$, then $N^t(i)=\emptyset$ for every $i$,
  so $\lceil|N^t(i)|/2\rceil=0$ and
  $|\{j\in N^t(i)\mid B_i^t\sim_{\tau_i}B_j^t\}|=0$.
  The inequality $0\le 0$ holds trivially, giving $S_i^{t+1}=1$.
  Note that this argument is independent of the bias functions $\beta_{i,j}$.
\end{proof}

\begin{corollary}[No perpetual global silence in \CBSOMminus]
\label{cor:no-perpetual-silence}
  In any \CBSOMminus\ model, for every $t\in\mathbb{N}$ there exists
  $i\in A$ such that $S_i^t=1$ or $S_i^{t+1}=1$.
\end{corollary}

The next result extends Lemma~\ref{lem:bounds} to \CBSOMminus.

\begin{restatable}[Bounded opinions under \CBSOMminus]{lemma}{cbsomBounds}
\label{lem:cbsom-bounds}
  Let $(G,\mathbf{B}^0,\mathbf{S}^0,\mu_G)$ be a \CBSOMminus\ model
  with continuous bias functions $\beta_{i,j}\in R$ for all edges.
  For any $t\in\mathbb{N}$ and all $i\in A$:
  \[\min(\mathbf{B}^t)\;\le\; B_i^{t+1}\;\le\;\max(\mathbf{B}^t).\]
\end{restatable}

\begin{proof}[Proof sketch]
  For any vocal neighbor $j\in N^t(i)$, let $x=B_j^t-B_i^t$.
  Since $\beta_{i,j}\in R$: $\operatorname{sign}(\beta_{i,j}(x))=
  \operatorname{sign}(x)$ and $|\beta_{i,j}(x)|\le|x|$, so the term
  $I_{ji}\cdot\beta_{i,j}(x)$ moves $B_i^t$ toward $B_j^t$ without
  overshooting.  The row-stochasticity of $I$ then yields the bound by a convexity argument.
  Full proof in Appendix~\ref{app:proofs}.
\end{proof}

\begin{corollary}[Monotone extremes in \CBSOMminus]
\label{cor:cbsom-mono}
  In any \CBSOMminus\ model with continuous $R$-biases,
  $\max(\mathbf{B}^t)$ is non-increasing and
  $\min(\mathbf{B}^t)$ is non-decreasing in $t$.
\end{corollary}

\begin{theorem}[Limits of extremes in \CBSOMminus]
\label{thm:cbsom-limits}
  For any \CBSOMminus\ model with continuous $R$-biases, there exist
  $U,L\in[0,1]$ with $L\le U$ such that
  $\lim_{t\to\infty}\max(\mathbf{B}^t)=U$ and
  $\lim_{t\to\infty}\min(\mathbf{B}^t)=L$.
\end{theorem}

\begin{proof}
  $\{\max(\mathbf{B}^t)\}$ is bounded below by $0$ and non-increasing
  (Corollary~\ref{cor:cbsom-mono});
  $\{\min(\mathbf{B}^t)\}$ is bounded above by $1$ and non-decreasing.
  Both limits exist by the Monotone Convergence
  Theorem~\cite{Houshang2014}.
\end{proof}

% ===================================================================
\FloatBarrier
\subsection{Scaled Linear Bias Functions}
\label{sec:scaled-linear}

\begin{definition}[Scaled linear bias family]
\label{def:scaled-linear}
  The \emph{scaled linear bias family} is
  \[\mathcal{F}_\alpha
    =\{f_\alpha:[-1,1]\to[-1,1]\mid f_\alpha(x)=\alpha x,\;
    \alpha\in(0,1)\}.\]
\end{definition}

The scalar $\alpha\in(0,1)$ encodes the degree of receptiveness:
as $\alpha\to 1$, $f_\alpha$ approaches the DeGroot identity; as
$\alpha\to 0$, it approaches the insular function.
Figure~\ref{fig:bias-functions-scaled-family} shows the family
alongside the DeGroot identity and the confirmation bias
function of~\cite{Alvim2024} on the domain $S=[-1,1]^2$.

\begin{figure}[!htbp]
  \centering%
  \includegraphics[width=0.54\textwidth]{bias_functions_scaled_family.pdf}
  \caption{Bias functions on $S=[-1,1]^2$.
    The identity $f(x)=x$ (DeGroot, green) marks the boundary between
    the Malleability region $M$ and the Receptive-Resistant region $R$.
    The confirmation bias $\mathrm{conf}(x)$ (blue,~\cite{Alvim2024})
    exhibits nonlinear resistance to distant opinions.
    Three members of $\mathcal{F}_\alpha$ are shown with
    $\alpha\in\{0.1,0.5,0.9\}$ (magenta, cyan, dark-red).
    In the first quadrant all scaled functions lie strictly below the
    identity, and in the third quadrant strictly above it, confirming
    membership in $R$}
  \label{fig:bias-functions-scaled-family}
\end{figure}

\begin{lemma}[Scaled linear biases lie in $R$]\label{lem:in-R}
  For every $\alpha\in(0,1)$, $f_\alpha\in R$.
\end{lemma}

\begin{proof}
  For $x>0$: $f_\alpha(x)=\alpha x\in(0,x)$, so $(x,f_\alpha(x))\in R$.
  For $x<0$: $f_\alpha(x)=\alpha x\in(x,0)$, so $(x,f_\alpha(x))\in R$.
  For $x=0$: $f_\alpha(0)=0$, so $(0,0)\in R$.
\end{proof}

By Lemma~\ref{lem:in-R}, all results of
Section~\ref{sec:analysis-structure} apply uniformly to
$\mathcal{F}_\alpha$.  Moreover, each $f_\alpha$ induces a
\emph{contraction} on the opinion spread
$R_t=\max(\mathbf{B}^t)-\min(\mathbf{B}^t)$.

\begin{restatable}[Geometric decay of opinion spread]{lemma}{geomDecay}
\label{lem:contraction}
  Let $(G,\mathbf{B}^0,\mathbf{S}^0,\mu_G)$ be a \CBSOMminus\ model
  on a clique with $n\ge 2$, with $\beta_{i,j}=f_{\alpha_{i,j}}$ for
  all edges. Let
  $\alpha_{\min}=\min_{(j,i)\in E}\alpha_{i,j}$ and
  $I_{\min}=\min_{(j,i)\in E}I_{ji}$.
  Then for every $\varepsilon>0$ there exists $t\in\mathbb{N}$ such
  that $R_t < \varepsilon$.
\end{restatable}

\begin{proof}[Proof sketch]
  Each round in which a vocal neighbor $j$ influences the max-agent $i$
  decreases the maximum by at least $I_{\min}\cdot\alpha_{\min}\cdot R_t$,
  since all terms in the update of the max-agent are non-positive.
  By Corollary~\ref{cor:no-perpetual-silence}, such events occur
  infinitely often.  The full proof (in Appendix~\ref{app:proofs})
  also handles the case where the min-agent is persistently silent,
  showing its opinion is pulled toward the maximum by vocal agents,
  again forcing $R_t\to 0$.
\end{proof}

% ===================================================================
\FloatBarrier
\subsection{Consensus in Clique Graphs}
\label{sec:consensus-cliques}

% -------------------------------------------------------------------
\subsubsection{Cliques with \texorpdfstring{$n\ge 3$}{n >= 3}}

\begin{restatable}[Consensus in \CBSOMminus\ cliques with $n\ge 3$]{theorem}{cbsomConsensus}
\label{thm:cbsom-consensus-n3}
  Let $(G,\mathbf{B}^0,\mathbf{S}^0,\mu_G)$ be a \CBSOMminus\ model
  on an $n$-agent clique with $n\ge 3$ and
  $\beta_{i,j}=f_{\alpha_{i,j}}\in\mathcal{F}_\alpha$ for all edges.
  Then all agents converge to consensus.
\end{restatable}

\begin{proof}[Proof sketch]
  By Theorem~\ref{thm:cbsom-limits}, limits $U$ and $L$ exist.
  Suppose $U-L>0$.  By Lemma~\ref{lem:contraction}, $R_t<U-L$ for
  some $t$; yet Corollary~\ref{cor:cbsom-mono} requires
  $R_t\ge U-L$ for all $t$: contradiction.  Hence $U=L$, and the
  Squeeze Theorem~\cite{Houshang2014} yields consensus.
  Full proof in Appendix~\ref{app:proofs}.
\end{proof}

\begin{remark}
  Theorem~\ref{thm:cbsom-consensus-n3} extends
  Aranda et al.'s consensus result for \SOMminus\
  cliques~\cite{aranda2024sound} (identity bias, $\alpha=1$) to the
  full scaled linear family $\mathcal{F}_\alpha$, and lifts
  Alvim et al.'s $R$-bias consensus
  result~\cite{Alvim2024} from the synchronous DeGroot setting to the
  Spiral-of-Silence setting.
\end{remark}

Figures~\ref{fig:clique3-conf} and~\ref{fig:clique3-scaled} illustrate
consensus convergence on a 3-agent clique under $\mathrm{conf}(x)$
and $f_\alpha$ biases, respectively.

\begin{figure}[!htbp]
  \centering
  \includegraphics[width=0.95\textwidth]{FourthExperimentClique3Confirmation.pdf}
  \caption{Consensus on a 3-agent clique under \CBSOMminus\
    with $\mathrm{conf}(x)$ bias on all edges
    (Theorem~\ref{thm:cbsom-consensus-n3}).
    All directed edges carry influence weight $0.5$;
    initial beliefs are $\mathbf{B}^0=(1,0.5,0)$; $\tau_i=1$}
  \label{fig:clique3-conf}
\end{figure}

\begin{figure}[!htbp]
  \centering
  \includegraphics[width=0.95\textwidth]{FifthExperimentScaledConfirmationCliqueC.pdf}
  \caption{Consensus on a 3-agent clique under \CBSOMminus\
    with scaled linear bias $f_\alpha$, $\alpha=0.5$
    (Theorem~\ref{thm:cbsom-consensus-n3}).
    Same parameters as Figure~\ref{fig:clique3-conf}.
    $f_\alpha$ converges faster than $\mathrm{conf}(x)$ due to
    stronger linear attenuation of disagreement at all scales}
  \label{fig:clique3-scaled}
\end{figure}

% -------------------------------------------------------------------
\subsubsection{The Dyadic Case \texorpdfstring{$n=2$}{n = 2}}

The \SOMminus\ model fails to converge on two-agent cliques,
exhibiting period-2 oscillations~\cite{aranda2024sound}.
Figure~\ref{fig:som-n2-oscillation} illustrates this pathology:
with symmetric influence weights $I_{12}=I_{21}=1.0$ and
$\tau_i=1$, agents permanently alternate their beliefs between $0$
and $1$ without reaching any agreement.

\begin{figure}[!htbp]
  \centering%
  \includegraphics[width=0.95\textwidth]{FirstExperimentN2.pdf}
  \caption{Non-convergence counterexample for \SOMminus\ on a dyadic
    clique ($n=2$).
    Symmetric influence weights $I_{12}=I_{21}=1.0$,
    $\mathbf{B}^0=(1,0)$, $\tau_i=1$.
    Agents permanently oscillate between $0$ and $1$, alternating
    speaking and silent roles; consensus is never reached.
    This is the pathology that Theorem~\ref{thm:cbsom-consensus-n2}
    resolves by introducing $\mathcal{F}_\alpha$-biases}
  \label{fig:som-n2-oscillation}
\end{figure}

We now show that $\mathcal{F}_\alpha$-biases eliminate this pathology.

\begin{lemma}[Persistent vocality for $n=2$]\label{lem:vocal-n2}
  Consider \CBSOMminus\ on two agents $\{1,2\}$ with $\tau_1=\tau_2=1$.
  Then $S_1^t=S_2^t=1$ for all $t\in\mathbb{N}$.
\end{lemma}

\begin{proof}
  With $\tau_i=1$, the condition $|B_i^t-B_j^t|\le 1$ holds for all
  $t$ since opinions lie in $[0,1]$.  Hence every agent always satisfies
  the majority criterion and remains vocal.
\end{proof}

\begin{theorem}[Consensus under scaled linear biases for $n=2$]
\label{thm:cbsom-consensus-n2}
  Consider \CBSOMminus\ on a two-agent clique $\{1,2\}$ with
  $I_{21},I_{12}>0$, biases $\beta_{i,j}(x)=\alpha_{i,j}x$ with
  $\alpha_{i,j}\in(0,1)$, and $\tau_1=\tau_2=1$.
  Then both agents converge to consensus.
\end{theorem}

\begin{proof}
  By Lemma~\ref{lem:vocal-n2}, $S_j^t=1$ for all $t$.
  Let $\lambda_{12}=I_{21}\alpha_{1,2}$ and
  $\lambda_{21}=I_{12}\alpha_{2,1}$.
  Eq.~\eqref{eq:cbsom-minus-b} reduces to:
  \begin{align*}
    B_1^{t+1}&=(1-\lambda_{12})B_1^t+\lambda_{12}B_2^t,\\
    B_2^{t+1}&=\lambda_{21}B_1^t+(1-\lambda_{21})B_2^t.
  \end{align*}
  Let $D^t=B_1^t-B_2^t$.  Subtracting:
  $D^{t+1}=\gamma D^t$ where
  $\gamma=1-\lambda_{12}-\lambda_{21}$.
  Since $\alpha_{i,j}\in(0,1)$ and $I_{ji}>0$ for both edges,
  $\lambda_{12}=I_{21}\alpha_{1,2}\in(0,I_{21})$ and
  $\lambda_{21}=I_{12}\alpha_{2,1}\in(0,I_{12})$ are both strictly
  positive.  Hence $\lambda_{12}+\lambda_{21}\in(0,\,I_{21}+I_{12})$.
  Since $I_{12},I_{21}\le 1$ and $\alpha_{i,j}<1$,
  $\lambda_{12}+\lambda_{21}<2$; combined with strict positivity,
  $\gamma=1-\lambda_{12}-\lambda_{21}\in(-1,1)$, so $|\gamma|<1$.
  Therefore $D^t=\gamma^t D^0\to 0$, and both agents converge to
  the same limit $v\in[0,1]$.
\end{proof}

\begin{remark}[Asymmetric bias case]\label{rem:n2-asymmetric}
  Theorem~\ref{thm:cbsom-consensus-n2} requires $\alpha_{i,j}\in(0,1)$
  for \emph{both} directed edges.  The asymmetric case---one edge with
  $\alpha=1$ (identity, no bias) and one with $\alpha\in(0,1)$---falls
  outside the stated hypotheses but is conjectured to converge; we
  leave a formal treatment for future work.
\end{remark}

\begin{remark}[Recovery of the $n=2$ pathology]
\label{rem:n2-recovery}
  The period-2 oscillation of \SOMminus~\cite{aranda2024sound}
  corresponds to $\gamma=-1$, requiring
  $\lambda_{12}+\lambda_{21}=2$.
  With $\alpha_{i,j}\in(0,1)$, $\lambda_{ij}=I_{ji}\alpha_{i,j}<1$
  strictly, so $|\gamma|<1$ always.
  In the symmetric case ($I_{12}=I_{21}=I$,
  $\alpha_{1,2}=\alpha_{2,1}=\alpha$),
  the contraction rate is $\gamma=1-2\alpha I$:
  larger $\alpha$ produces oscillatory (but still convergent) dynamics,
  while smaller $\alpha$ yields monotone smooth convergence.
\end{remark}

Figures~\ref{fig:n2-alpha01} and~\ref{fig:n2-alpha09} illustrate how
$\alpha$ controls the qualitative character of convergence in
Case~2 of Theorem~\ref{thm:cbsom-consensus-n2}.

\begin{figure}[!htbp]
  \centering
  \includegraphics[width=0.95\textwidth]{2alpha01.pdf}
  \caption{Consensus on a two-agent clique under \CBSOMminus\ with
    scaled linear bias $f_\alpha$, $\alpha=0.1$
    (Case~2 of Theorem~\ref{thm:cbsom-consensus-n2}).
    Parameters: $I_{12}=I_{21}=1.0$, $\mathbf{B}^0=(1,0)$, $\tau_i=1$.
    With $\gamma=1-2\alpha I=0.8$ the opinion gap shrinks monotonically
    (``inverse trumpet'' profile)}
  \label{fig:n2-alpha01}
\end{figure}

\begin{figure}[!htbp]
  \centering
  \includegraphics[width=0.95\textwidth]{2alpha09.pdf}
  \caption{Consensus on a two-agent clique under \CBSOMminus\ with
    scaled linear bias $f_\alpha$, $\alpha=0.9$.
    Same parameters as Figure~\ref{fig:n2-alpha01}.
    With $\gamma=-0.8$ the gap contracts oscillatorily but with
    strictly decreasing amplitude, reminiscent of the \SOMminus\
    pathology (Figure~\ref{fig:som-n2-oscillation}) but convergent.
    The special case $\alpha=0.5$ reaches consensus
    exactly at round~1: $B_1^1=B_2^1=0.5$}
  \label{fig:n2-alpha09}
\end{figure}

% ===================================================================
\FloatBarrier
\subsection{Summary}
\label{sec:analysis-summary}

Table~\ref{tab:consensus-summary} summarizes the consensus landscape
across models and clique sizes.

\begin{table}[ht]
\centering%
\caption{Consensus in clique graphs: \SOMminus\ vs.\ \CBSOMminus\
  with $\mathcal{F}_\alpha$-biases}
\label{tab:consensus-summary}
\begin{tabular}{lcc}
  \hline
  \textbf{Model} & \textbf{$n=2$} & \textbf{$n\ge 3$}\\
  \hline
  \SOMminus~\cite{aranda2024sound}
    & oscillates & converges\\
  \CBSOMminus\ $+$ $\mathcal{F}_\alpha$ \emph{(this work)}
    & converges & converges\\
  \hline
\end{tabular}
\end{table}

The decoupling between bias and silence---formalized in
Lemma~\ref{lem:no-perpetual-silence}---is the structural key: silence
determines \emph{who speaks} while bias determines \emph{how the
spoken opinion is absorbed}.  Consensus emerges from the joint action
of the contraction induced by $\mathcal{F}_\alpha$
(Lemma~\ref{lem:contraction}) and the impossibility of perpetual
silence (Corollary~\ref{cor:no-perpetual-silence}).
